\documentclass[10pt,a4paper]{amsart}
\usepackage{amsmath,amssymb,amsthm,graphicx}
% Retained preamble for the Stacks Project chapter divisors.tex (Divisors).
% The chapter file pinned-original-source.tex carries only \input{preamble} on its first line, so no
% preamble text lives inside the chapter itself.  The Stacks Project uses one single shared preamble for
% every chapter; the snapshot stages it verbatim in sources/stacks-chow/preamble.tex.  The block below is
% copied mechanically, line for line, from that file's lines 14-18 (the Xy-pic setup) and lines 154-217
% (the native theorem declarations together with the mathematical macros).  Nothing was added, removed or
% redefined.  The wrapper class and the amsmath/amssymb/amsthm/graphicx base packages are supplied by the
% assembler, so the class/\externaldocument part of the shared preamble is not repeated here.
% For commutative diagrams we use Xy-pic
\usepackage[all]{xy}

% We use 2cell for 2-commutative diagrams.
\xyoption{2cell}
% Theorem environments.
%
\theoremstyle{plain}
\newtheorem{theorem}[subsection]{Theorem}
\newtheorem{proposition}[subsection]{Proposition}
\newtheorem{lemma}[subsection]{Lemma}

\theoremstyle{definition}
\newtheorem{definition}[subsection]{Definition}
\newtheorem{example}[subsection]{Example}
\newtheorem{exercise}[subsection]{Exercise}
\newtheorem{situation}[subsection]{Situation}

\theoremstyle{remark}
\newtheorem{remark}[subsection]{Remark}
\newtheorem{remarks}[subsection]{Remarks}

\numberwithin{equation}{subsection}

% Macros
%
\def\lim{\mathop{\mathrm{lim}}\nolimits}
\def\colim{\mathop{\mathrm{colim}}\nolimits}
\def\Spec{\mathop{\mathrm{Spec}}}
\def\Hom{\mathop{\mathrm{Hom}}\nolimits}
\def\Ext{\mathop{\mathrm{Ext}}\nolimits}
\def\SheafHom{\mathop{\mathcal{H}\!\mathit{om}}\nolimits}
\def\SheafExt{\mathop{\mathcal{E}\!\mathit{xt}}\nolimits}
\def\Sch{\mathit{Sch}}
\def\Mor{\mathop{\mathrm{Mor}}\nolimits}
\def\Ob{\mathop{\mathrm{Ob}}\nolimits}
\def\Sh{\mathop{\mathit{Sh}}\nolimits}
\def\NL{\mathop{N\!L}\nolimits}
\def\CH{\mathop{\mathrm{CH}}\nolimits}
\def\proetale{{pro\text{-}\acute{e}tale}}
\def\etale{{\acute{e}tale}}
\def\QCoh{\mathit{QCoh}}
\def\Ker{\mathop{\mathrm{Ker}}}
\def\Im{\mathop{\mathrm{Im}}}
\def\Coker{\mathop{\mathrm{Coker}}}
\def\Coim{\mathop{\mathrm{Coim}}}

% Boxtimes
%
\DeclareMathSymbol{\boxtimes}{\mathbin}{AMSa}{"02}

%
% Macros for moduli stacks/spaces
%
\def\QCohstack{\mathcal{QC}\!\mathit{oh}}
\def\Cohstack{\mathcal{C}\!\mathit{oh}}
\def\Spacesstack{\mathcal{S}\!\mathit{paces}}
\def\Quotfunctor{\mathrm{Quot}}
\def\Hilbfunctor{\mathrm{Hilb}}
\def\Curvesstack{\mathcal{C}\!\mathit{urves}}
\def\Polarizedstack{\mathcal{P}\!\mathit{olarized}}
\def\Complexesstack{\mathcal{C}\!\mathit{omplexes}}
% \Pic is the operator that assigns to X its picard group, usage \Pic(X)
% \Picardstack_{X/B} denotes the Picard stack of X over B
% \Picardfunctor_{X/B} denotes the Picard functor of X over B
\def\Pic{\mathop{\mathrm{Pic}}\nolimits}
\def\Picardstack{\mathcal{P}\!\mathit{ic}}
\def\Picardfunctor{\mathrm{Pic}}
\def\Deformationcategory{\mathcal{D}\!\mathit{ef}}

% --- appended 2026-10-12 (Chief): the project own definitions of the reference/slogan/history
% display environments, copied mechanically from the shared Stacks preamble lines 8-12
% (\usepackage{verbatim} plus the three \newenvironment lines, which the chapter preamble keeps
% verbatim from the project preamble). Nothing added, removed or redefined: these environments
% carry no mathematical content and the project itself typesets them as comments. Without these
% three lines a counted statement whose own text uses them fails with Environment slogan undefined.
\usepackage{verbatim}
\newenvironment{reference}{\comment}{\endcomment}
\newenvironment{slogan}{\comment}{\endcomment}
\newenvironment{history}{\comment}{\endcomment}
\usepackage{cleveref}
\crefname{originallemma}{original lemma}{original lemmas}
\Crefname{originallemma}{Original lemma}{Original lemmas}
\makeatletter
\newlabel{lemma-meromorphic-functions-integral-scheme}{{01X5}{source}{}{sourcepointer-1}{}}
\newlabel{lemma-meromorphic-functions-integral-scheme@cref}{{[originallemma][0][]01X5}{[1][0][]source}{}{}{}}
\makeatother
\crefname{originallemma}{original lemma}{original lemmas}
\Crefname{originallemma}{Original lemma}{Original lemmas}
\makeatletter
\newlabel{lemma-reduced-finite-irreducible}{{02OW}{source}{}{sourcepointer-2}{}}
\newlabel{lemma-reduced-finite-irreducible@cref}{{[originallemma][0][]02OW}{[1][0][]source}{}{}{}}
\makeatother
\crefname{originallemma}{original lemma}{original lemmas}
\Crefname{originallemma}{Original lemma}{Original lemmas}
\makeatletter
\newlabel{lemma-regular-meromorphic-section-exists-noetherian}{{0EMI}{source}{}{sourcepointer-3}{}}
\newlabel{lemma-regular-meromorphic-section-exists-noetherian@cref}{{[originallemma][0][]0EMI}{[1][0][]source}{}{}{}}
\makeatother
\title{Invertible modules have regular meromorphic sections in three cases}
\author{The Stacks Project Authors}
\date{stacks-project:divisors.tex:local-snapshot-20261012; GNU FDL 1.2 or later}
\begin{document}
\maketitle
\setcounter{subsection}{1}
\noindent\textit{Editorial note.} All delivered mathematical text is the author's own, copied from the pinned Stacks Project chapter file divisors.tex; no statement or proof step was rewritten, reordered or removed.\par\medskip
\noindent\textit{Editorial note.} Carrier window 3: statement and proof both lie inside the window 6000-8000 granted to this lane by the owner lane lane-s13 on 2026-10-14.\par\medskip
\noindent\textit{Editorial note.} Cases (1) and (2) are argued in place; case (3) is a single named input, and that shape is recorded here for the reader.\par\medskip
\noindent\textit{Editorial note.} Preamble: sources/stacks-divisors/retained-preamble.tex (88 lines, verbatim present).\par\medskip
\noindent\textit{Allowed standard prerequisites (editorial).}\begin{itemize}
\item Meromorphic functions and sections on an integral scheme. Lemma lemma-meromorphic-functions-integral-scheme, used at proof lines 6524-6528. Identifies the two sheaves with constant sheaves at the generic point.
\item Description of K\_X and K\_X(L) for a reduced scheme with finitely many components. Lemma lemma-reduced-finite-irreducible (6424-6456, produced as S15-div-reduced-finite-irreducible), used at proof lines 6534-6537. Makes the product of the chosen nonzero elements a meromorphic section.
\item Existence of a regular meromorphic section in the Noetherian case. Lemma lemma-regular-meromorphic-section-exists-noetherian, used at proof line 6541. Is case (3), cited by name.
\end{itemize}\medskip
\ifdefined\hypertarget\hypertarget{sourcepointer-1}{}\fi
\noindent\textit{Source reference.} Original lemma 01X5: focused-seeds/AGP/sources/stacks-snapshot-20261011/tags.txt, source lines 2506--2506. Navigation pointer into the Stacks Project; the standard fact it supplies is applied inside the counted proof span and the argument resting on it stays there.\par\medskip
\ifdefined\hypertarget\hypertarget{sourcepointer-2}{}\fi
\noindent\textit{Source reference.} Original lemma 02OW: focused-seeds/AGP/sources/stacks-snapshot-20261011/tags.txt, source lines 3531--3531. Navigation pointer into the Stacks Project; the standard fact it supplies is applied inside the counted proof span and the argument resting on it stays there.\par\medskip
\ifdefined\hypertarget\hypertarget{sourcepointer-3}{}\fi
\noindent\textit{Source reference.} Original lemma 0EMI: focused-seeds/AGP/sources/stacks-snapshot-20261011/tags.txt, source lines 18308--18308. Navigation pointer into the Stacks Project; the standard fact it supplies is applied inside the counted proof span and the argument resting on it stays there.\par\medskip
\begin{lemma}
\label{lemma-regular-meromorphic-section-exists}
Let $X$ be a scheme.
Let $\mathcal{L}$ be an invertible $\mathcal{O}_X$-module.
In each of the following cases $\mathcal{L}$ has a regular meromorphic
section:
\begin{enumerate}
\item $X$ is integral,
\item $X$ is reduced and any quasi-compact open has a finite
number of irreducible components,
\item $X$ is locally Noetherian and has no embedded points.
\end{enumerate}
\end{lemma}

\begin{proof}
In case (1) let $\eta \in X$ be the generic point. We have seen in
Lemma \ref{lemma-meromorphic-functions-integral-scheme}
that $\mathcal{K}_X$, resp.\ $\mathcal{K}_X(\mathcal{L})$
is the constant sheaf with value
$\kappa(\eta)$, resp.\ $\mathcal{L}_\eta$.
Since $\dim_{\kappa(\eta)} \mathcal{L}_\eta = 1$
we can pick a nonzero element $s \in \mathcal{L}_\eta$.
Clearly $s$ is a regular meromorphic section of $\mathcal{L}$.
In case (2) pick $s_\eta \in \mathcal{L}_\eta$ nonzero
for all generic points $\eta$ of $X$; this is possible
as $\mathcal{L}_\eta$ is a $1$-dimensional vector space
over $\kappa(\eta)$.
It follows immediately from the description of
$\mathcal{K}_X$ and $\mathcal{K}_X(\mathcal{L})$
in Lemma \ref{lemma-reduced-finite-irreducible}
that $s = \prod s_\eta$ is a regular meromorphic section of $\mathcal{L}$.
Case (3) is Lemma \ref{lemma-regular-meromorphic-section-exists-noetherian}.
\end{proof}

\end{document}
